% Generated draft: requires mathematical and source review.
\documentclass[11pt]{article}
\usepackage[margin=25mm]{geometry}
\usepackage{fontspec}
\setmainfont{Latin Modern Roman}
\usepackage{amsmath,amssymb,mathtools,unicode-math}
\setmathfont{Latin Modern Math}
\usepackage{enumitem,needspace,xurl,hyperref}
\setlength{\parindent}{0pt}
\setlength{\parskip}{5pt}
\setlength{\emergencystretch}{3em}
\newcommand{\N}{\mathbb{N}}
\newcommand{\Z}{\mathbb{Z}}
\newcommand{\Q}{\mathbb{Q}}
\newcommand{\R}{\mathbb{R}}
\newcommand{\C}{\mathbb{C}}
\newcommand{\F}{\mathbb{F}}
\newcommand{\A}{\mathbb{A}}
\newcommand{\PP}{\mathbb{P}}
\newcommand{\E}{\mathbb{E}}
\newcommand{\eps}{\varepsilon}
\newcommand{\dd}{\,\mathrm{d}}
\begin{document}
\section*{UnsolvedMath ID: 11000328}
\textbf{Source problem code:} catalogue entry
\subsection{Definitions and mathematical statement}
Let $\Sigma_g$ be a closed oriented genus-$g$ surface. Write $\operatorname{Diff}_+^\delta(\Sigma_g)$ for its orientation-preserving diffeomorphism group endowed with the discrete topology, and $\operatorname{Symp}^\delta(\Sigma_g,\omega_g)$ for the discrete group preserving a specified area form $\omega_g$. The superscript $\delta$ is essential. For a discrete group $G$, $BG$ is its classifying space, so $H_i(BG;A)=H_i(G;A)$ for the chosen coefficient system $A$.

\textbf{Problem 12.2.} Study whether the homology groups
\[
 H_i(B\operatorname{Diff}_+^\delta(\Sigma_g);A)
\]
stabilize as $g$ increases. Ask the same question for
\[
 H_i(B\operatorname{Symp}^\delta(\Sigma_g,\omega_g);A).
\]
The usual stability target is, for each fixed degree $i$, an eventual isomorphism under specified genus-increasing comparison maps, with a threshold allowed to depend on $i$. The coefficient convention and comparison maps are not provided in the excerpt. Genus-increasing extension maps are naturally defined for groups fixing a boundary collar or a disk; using those versions to interpret the closed-surface question requires a stated capping/comparison argument. It is not legitimate to assume every diffeomorphism of a closed surface has a chosen extension over a connected sum. For the symplectic version, the area forms and compatibility of any gluing maps also require specification.

The question is not ordinary mapping-class-group homological stability. Passing from discrete diffeomorphisms or symplectomorphisms to mapping classes forgets substantial group data. The supplied summary invokes Harer and Madsen--Weiss to label the question solved, but gives no result identifying their mapping-class stability statements with these discrete-group homology groups. That reasoning therefore does not verify the source problem. No present positive or negative stability result is independently asserted here.
\subsection{Short English statement}
Ask whether genus stabilization holds for the homology of discrete surface diffeomorphism groups and discrete area-preserving groups. The discrete topology, coefficients, and genus comparison maps must be kept distinct from mapping-class-group stability.
\subsection{Sources}
\begin{itemize}
\item Local UnsolvedMath registry, permanent ID 11000328, source code catalogue entry; no separate proposer or proposal year supplied.
\item Farb, edited volume \emph{Problems on Mapping Class Groups and Related Topics}, chapter ``Cohomological structure of the mapping class group and beyond,'' Problem 12.2, PDF page 380, in text bearing S. Morita's attribution; \url{https://www.math.uchicago.edu/~farb/papers/mcgbook.pdf}. The registry records PDF-font-structure extraction and access on 2026-07-29; the inventory distinguishes 329 labeled headings by source page.
\item The research summary names Harer and Madsen--Weiss but provides no precise primary theorem for the discrete groups in this problem.
\item The source's closing acknowledgments credit R. Hain, N. Kawazumi, D. Kotschick, M. Matsumoto, H. Nakamura, and T. Sakasai for discussions and information; these are acknowledgments, not attributions of a stability solution.
\end{itemize}
\end{document}
